\section{Discussion}

\subsection{Paper-to-code is not paper-to-engineering}
A generated function can be syntactically correct and numerically stable while still using a misread table cell, an ambiguous unit, or an unsupported convention. For engineering reuse, source provenance and discrepancy handling are therefore part of the computational method rather than documentation added after the fact.

\subsection{Human review as a controlled state transition}
Automation classifies and escalates discrepancies, but does not decide engineering acceptability. Human decisions are recorded append-only using dispositions such as resolved, bounded, accepted with rationale, or deferred. An unresolved high-priority discrepancy can block study promotion while leaving the original source evidence unchanged. These decisions do not grant engineering qualification.

\subsection{Evidence graphs}
The evidence graph links source identity, selected targets, callable methods, comparisons, discrepancies, automated assessments, and human decisions. Graph synchronization is designed to be additive and idempotent. Coverage is a provenance property only; a complete graph does not imply that the engineering method is verified or qualified.

\subsection{Implications for research software}
The source-backed mismatches in the two development cases suggest that an evidence-aware system may contribute value even when it cannot fully reproduce a publication. Detecting that a source is internally inconsistent, or that a result cannot be independently reconstructed without additional assumptions, can be more useful for engineering review than forcing an apparent numerical match.
